\title{FlashAttention-3: Fast and Accurate Attention with Asynchrony and Low-precision}

\begin{abstract}
The attention mechanism serves as a fundamental component within the widely adopted Transformer framework, yet it frequently limits the performance of large language models, especially for tasks requiring extended contexts.
\textsc{FlashAttention} introduced a strategy that improves attention computation speed on GPUs by reducing unnecessary memory access. Nonetheless, this method does not fully harness advancements available in modern hardware; for instance, \textsc{FlashAttention-2} utilizes only 35\% of the H100 GPU's capabilities. In this work, we introduce three key strategies to enhance attention throughput on Hopper GPUs: by utilizing the asynchronous features of the Tensor Cores and TMA, we (1) achieve overlap between computation and memory transfer through warp-specialization, (2) alternate block-wise matrix multiplications with softmax computations, and (3) apply block-level quantization and incoherent processing, thereby making use of FP8 low-precision supported by hardware. Our proposed technique, \textsc{FlashAttention-3}, provides a 1.5-2.0$\times$ increase in speed on H100 GPUs, reaching a throughput of 740 TFLOPs/s in FP16 mode (75\% hardware utilization), and nearly 1.2 PFLOPs/s using FP8. Further, we show that FP8 \textsc{FlashAttention-3} reduces numerical error by 2.6$\times$ relative to a standard FP8 attention baseline.
\end{abstract}

\section{Introduction}
\label{sec:intro}

Within the Transformer architecture~\citep{vaswani2017attention}, the chief computational challenge arises from the attention mechanism, primarily due to the quadratic complexity incurred when evaluating self-attention scores across queries and keys as sequence lengths increase. Enabling the attention mechanism to efficiently process extended sequences paves the way for advances in a variety of domains: from modeling and reasoning over several lengthy documents~\citep{guo2021longt5,shaham2022scrolls,peng2023yarn} and large-scale codebases~\citep{roziere2023code, li2023starcoder}, to handling new data types such as high-resolution imagery~\citep{chen2022scaling}, audio signals~\citep{gulati2020conformer}, and video streams~\citep{ho2022video}. This also broadens practical applications, allowing for enhanced user experiences with extended interaction histories~\citep{sun2019bert4rec} and improved agent planning across prolonged workflows~\citep{yao2022react}. These factors have fueled considerable research toward accelerating attention for longer contexts, using approaches like approximation methods~\citep{katharopoulos2020transformers,choromanski2020rethinking, tay2020efficient}, software-level optimizations~\citep{rabe2021self,dao2022flashattention,kwon2023efficient}, or exploring novel architectures~\citep{peng2023rwkv,sun2023retentive,gu2023mamba}.

This study extends the advancements made by \citet{dao2022flashattention}, who formulated precise attention algorithms by embedding the GPU's architectural properties and execution behavior directly into their design paradigm. In \citep{dao2022flashattention}, Dao et al. unveiled \textsc{FlashAttention}, an innovative tiling scheme for parallelizing attention that dispenses with unnecessary accesses to slow global memory by consolidating all attention operations within a unified GPU kernel. Subsequently, \citet{dao2023flashattention2} introduced \textsc{FlashAttention-2}, restructuring the algorithm to also leverage parallelism across the sequence length dimension and to carry out the inner loop of the forward pass over partitions of the key and value matrices. This redesign enhanced workload distribution and GPU utilization. Despite these improvements, our evaluation indicates that \textsc{FlashAttention-2} attains only limited efficiency on contemporary GPU hardware, particularly when compared to high-performance matrix-multiplication (GEMM) kernels—reporting figures like 35\% versus 80-90\% on the Hopper H100 GPU. Some of this gap is due to low-level implementation distinctions, notably the omission of Hopper-tailored instructions in favor of those intended for Ampere architecture when interfacing with Tensor Cores. Recent contributions such as ThunkerKitten~\citep{spector2024thunder} and cuDNN 9~\citep{cudnn9} demonstrate that utilizing Hopper-centric instructions alongside tiled abstractions significantly accelerates attention computation and reduces implementation complexity.

At a deeper level, the algorithm behind \textsc{FlashAttention-2} operates within a basic synchronous framework, and does not explicitly integrate asynchronous execution or low-precision arithmetic into its approach. Asynchrony emerges as a consequence of specialized hardware components developed to enhance the performance of critical operations in machine learning workloads, such as dedicated units for matrix multiplications (Tensor Cores) and for memory access (Tensor Memory Accelerator -- TMA), which function independently from the main CUDA cores responsible for logical operations and general arithmetic tasks. The use of reduced numerical precision formats, for instance, FP8 with Hopper or FP4 with Blackwell, continues an established strategy—building on prior implementations like FP16 (Pascal, 2017) and BF16 (Ampere, 2020)—where lower bit-width enables a twofold or fourfold increase in computational throughput while maintaining constant power and die size. The capabilities of Hopper relevant to these advancements are discussed in detail in \cref{subsec:hardware}.

The central technical obstacle lies in re-engineering \textsc{FlashAttention-2} to leverage these hardware innovations. Efficient use of asynchrony demands computational overlap between matrix multiplication and softmax calculations, despite the dependency structure between them, whereas employing low-precision necessitates strategies for reducing quantization artifacts, with particular attention to handling outlier activations typical in large language models~\citep{dettmers2208llm, sun2024massive}.

In pursuit of this goal, we introduce \textsc{FlashAttention-3}, presenting and integrating three novel techniques that advance performance on modern GPU architectures.\footnote{Our discussion focuses on NVIDIA's Hopper architecture, but the proposed algorithm is applicable to any GPU platform that supports strong asynchronous execution and efficient low-precision operations.}

\begin{enumerate}

\item \textbf{Producer-Consumer asynchrony:} We introduce a warp-specific software pipelining approach designed to leverage the decoupling of data transfer and computation on Tensor Cores. This is achieved by assigning data producers and consumers to distinct warps, thereby broadening the scope for masking both memory access delays and instruction issuing latency within the algorithm.

\item \textbf{Hiding softmax under asynchronous block-wise GEMMs:} The non-GEMM operations required for softmax, including floating-point multiply-add and exponentiation, have lower throughput and are overlapped with the asynchronous execution of WGMMA instructions for GEMM. To facilitate this, we refactor the \textsc{FlashAttention-2} method to eliminate key sequential dependencies between softmax and GEMM phases. In particular, in our two-stage method, softmax operates on one block of the scores matrix concurrently with WGMMA, which asynchronously calculates the subsequent block.

\item \textbf{Hardware-accelerated low-precision GEMM:} The forward pass algorithm is modified to harness FP8 Tensor Cores for GEMM operations, resulting in approximately a twofold increase in empirical TFLOPs/s. Accommodating this necessitates reconciling WGMMA's memory layout expectations regarding FP32 accumulators and FP8 operands. To address the potential precision degradation inherent to FP8 computation, we employ strategies such as block quantization and incoherent processing.
\end{enumerate}

To empirically assess the effectiveness of our approach, we evaluate \textsc{FlashAttention-3} on the H100 SXM5 GPU across multiple settings. Our results demonstrate that: (1) the FP16 configuration delivers a 1.5-2.0$\times$ performance boost compared to \textsc{FlashAttention-2} during the forward pass (with throughput peaking at 740 TFLOPs/s), and achieves a 1.5-1.75$\times$ acceleration in the backward pass; (2) FP8 attains nearly 1.2 PFLOPs/s; and (3) for longer sequence lengths, FP16 surpasses other methods, while FP8 maintains competitive performance\footnote{Specifically, with a head dimension of 64, FP8 in \textsc{FlashAttention-3} leads, whereas for head dimensions 128 and 256, FP8 matches performance when causal masking is absent but lags when masking is present.} compared to NVIDIA’s cuDNN attention implementation. Additionally, our findings confirm that FP16 in \textsc{FlashAttention-3} produces equivalent numerical accuracy to \textsc{FlashAttention-2}, and outperforms standard attention mechanisms due to storing intermediate steps (such as softmax rescaling) in FP32. Furthermore, FP8 mode in \textsc{FlashAttention-3}, utilizing block quantization with incoherent processing, achieves 2.6$\times$ greater accuracy than conventional attention employing per-tensor quantization, particularly in situations where outlier features are present.

We have released \textsc{FlashAttention-3} under a permissive license\footnote{\textsc{FlashAttention-3}
  is available at \texttt{https://github.com/Dao-AILab/flash-attention}}, and intend to collaborate with PyTorch and Hugging Face libraries to maximize its accessibility for researchers and developers.

\section{Background: Multi-Head Attention and GPU Characteristics}
\label{sec:background}

\subsection{Multi-Head Attention}
\label{subsec:multi_head_attn}

Let $\mathbf{Q}, \mathbf{K}, \mathbf{V} \in \mathbb{R}^{N \times d}$ be the query, key and value input sequences associated to a single head, where $N$ is the sequence length and $d$ is the head dimension. Then the attention output $\mathbf{O}$ is computed as:

\begin{equation*}
  \mathbf{S} = \alpha \mathbf{Q} \mathbf{K}^\top \in \mathbb{R}^{N \times N}, \quad \mathbf{P} = \mathrm{softmax}(\mathbf{S}) \in \mathbb{R}^{N \times N}, \quad \mathbf{O} = \mathbf{P}\mathbf{V} \in \mathbb{R}^{N \times d},
\end{equation*}

In this context, the $\mathrm{softmax}$ function is computed for each row, and commonly $\alpha$ is chosen as $1/\sqrt{d}$ to scale the values appropriately. To maintain numerical stability when working with exponentials, it is standard to subtract $\mathrm{rowmax}(\mathbf{S})$ from $\mathbf{S}$. When employing multi-head attention (MHA), distinct query, key, and value projection parameters are used for each head. These operations are executed concurrently across different heads and batches, resulting in the assembled output tensor.

Now let $\phi$ be a scalar loss function and let $\mathbf{d}(-) = \partial \phi / \partial (-)$ be notation for the gradient. Given the output gradient $\mathbf{dO} \in \mathbb{R}^{N \times d}$, we compute $\mathbf{dQ}$, $\mathbf{dK}$, and $\mathbf{dV}$ according to the chain rule as follows:

\begin{align*}
  \mathbf{dV} &= \mathbf{P}^\top \mathbf{dO} \in \mathbb{R}^{N \times d} \\
  \mathbf{dP} &= \mathbf{dO} \mathbf{V}^\top \in \mathbb{R}^{N \times N} \\
  \mathbf{dS} &= \mathrm{dsoftmax} (\mathbf{dP}) \in \mathbb{R}^{N \times N} \\
  \mathbf{dQ} &= \alpha \mathbf{dS} \mathbf{K} \in \mathbb{R}^{N \times d} \\
  \mathbf{dK} &= \alpha \mathbf{dS}^\top \mathbf{Q} \in \mathbb{R}^{N \times d},
\end{align*}

In this context, the relationship $\mathbf{d}s = (\mathrm{diag}(p) - p p^\top)\mathbf{d}p$ holds when $p = \mathrm{softmax}(s)$, viewing $p$ as a function dependent on the vector $s$. We denote the row-wise application of this expression by $\mathrm{dsoftmax}(\mathbf{dP})$. Moreover, this calculation can be efficiently executed in parallel across all heads and batches during the backward propagation phase of MHA.

\subsection{GPU hardware characteristics and execution model}
\label{subsec:hardware}

We outline the elements of the GPU execution model that are pertinent to \textsc{FlashAttention-3}, emphasizing the NVIDIA Hopper architecture as a specific realization of these principles.

\paragraph{Memory hierarchy:} The memory system on the GPU is structured as a tiered set of data storage regions, where greater bandwidth is generally associated with reduced capacity (\cref{tab:gpu-hierarchy})\footnote{\citet{luo2024benchmarking} cites a shared memory bandwidth of 128 bytes per clock cycle for each SM, which we scale by 132 SMs and a boost frequency of 1830 MHz.}. At the largest and slowest level is global memory (GMEM), also known as HBM, which is implemented with off-chip DRAM and is universally accessible by all streaming multiprocessors (SMs). Information residing in GMEM is automatically loaded into the on-chip L2 cache. Closer to the core, each SM possesses a compact, on-chip, software-controlled and extensively banked buffer termed shared memory (SMEM). The innermost storage consists of the register file located within each SM.

\paragraph{Thread hierarchy:} In the GPU programming framework, threads—individual units of execution—are organized into several hierarchical groupings. Starting from the smallest, threads are bundled into warps, each consisting of 32 threads; these warps are then aggregated into warpgroups, typically containing 4 adjacent warps. Above this, groups of threads operate together as threadblocks (also referred to as cooperative thread arrays or CTAs). In some architectures such as Hopper, multiple threadblocks can further combine to form threadblock clusters. The largest collection at the top level is known as a grid.

The relationship between these two hierarchies is highly interconnected. Threads belonging to a single CTA are simultaneously assigned to the same SM, while CTAs located within a shared cluster are dispatched together onto the same GPC. All threads within a CTA can access SMEM directly, but each thread possesses up to 256 individual registers (RMEM) that are exclusive to its own execution context.

\paragraph{Asynchrony and warp-specialization:}

GPUs function as throughput-oriented processors, leveraging both concurrent execution and asynchronous operations to mask delays associated with memory access and processing. In the context of asynchronous data transfers between GMEM and SMEM, Hopper introduces the Tensor Memory Accelerator (TMA), a specialized hardware component \cite[\S7.29]{cuda}. Additionally, setting itself apart from earlier designs like Ampere, Hopper’s Tensor Core—accessible through the warpgroup-scoped WGMMA instruction \cite[\S9.7.14]{ptx}—executes asynchronously and is capable of directly reading its inputs from shared memory.

The inclusion of hardware mechanisms for asynchrony enables the implementation of warp-specialized kernels, wherein the warps within a CTA assume dedicated producer or consumer functions, responsible exclusively for either computation or data transfers. This design, in general, enhances the compiler’s capacity to produce highly efficient instruction schedules \citep{warp-specialization-2011}. Moreover, Hopper introduces the ability to dynamically assign register resources among warpgroups through \verb|setmaxnreg| \cite[\S9.7.17.1]{ptx}, ensuring that warps engaged in MMAs receive an increased allocation of RMEM compared to those primarily handling TMA operations, which may require as little as a single thread.

\paragraph{Low-precision number formats:}
\label{sec:low-precision-gpu}
Current-generation GPUs incorporate dedicated hardware modules designed to enhance computations with reduced precision. For instance, the WGMMA operation can utilize FP8 Tensor Cores on the Hopper architecture, achieving double the throughput per streaming multiprocessor (SM) relative to operations using FP16 or BF16 formats.

Properly using FP8 WGMMA requires awareness of the operand layout restrictions. When performing a GEMM operation that computes $A \times B^{\top}$ for matrices $A$ ($M\times K$) and $B$ ($N\times K$), the term \emph{mn-major} denotes an operand that is contiguous in either the $M$ or $N$ (outer) dimension, while \emph{k-major} describes contiguity in the $K$ (inner) dimension. For FP16 WGMMA, SMEM operands may be provided in either mn-major or k-major format. Conversely, FP8 WGMMA only permits operands in k-major layout. Additionally, in cases such as attention mechanisms where consecutive GEMMs are fused into a single kernel, mismatched layouts between FP32 accumulators and FP8 operands can impede the execution of subsequent FP8 WGMMAs.

Within the framework of attention mechanisms, these constraints on layout require specific alterations to the structure of an FP8 algorithm, as detailed in \cref{sec:algofp8}.

\subsection{Standard Attention and Flash Attention} In line with the conventions set by \citet{dao2022flashattention}, we refer to \textbf{standard attention} as a GPU-based approach that explicitly stores the intermediate matrices $\mathbf{S}$ and $\mathbf{P}$ in HBM. The principal innovation behind \textsc{FlashAttention} was to introduce a local softmax reduction, thereby minimizing costly intermediate memory operations and integrating attention computation within a unified kernel. This local softmax operation is implemented in lines \ref{code-ws:softmax_start}-\ref{code-ws:softmax_end} of the consumer mainloop described in \cref{alg:flash3_wgmma_ws_only}, along with blockwise normalization of $\mathbf{O}$. A concise derivation confirming that this approach correctly calculates $\mathbf{O}$ can be found in \cite[\S 2.3.1]{dao2023flashattention2}.

\section{FlashAttention-3: Algorithm}
\label{sec:algo}

This section outlines the \textsc{FlashAttention-3} algorithm. To streamline our discussion, we center our attention on the forward pass; the backward pass procedure is addressed in~\cref{sec:algo_ws_bwd}. Our approach begins by illustrating the incorporation of warp-specialization and a circular SMEM buffer within the foundational structure of \textsc{FlashAttention-2}. Subsequently, we elaborate on leveraging WGMMA asynchrony to establish a GEMM-softmax two-phase pipeline that enables computation and memory overlap. Lastly, we delineate the necessary adjustments for FP8, covering both format consistency as well as accuracy enhancements, which are achieved using block quantization and handling of incoherent operations.

\subsection{Producer-Consumer asynchrony through warp-specialization and pingpong scheduling}
\label{sec:algo_ws}

\paragraph{Warp-specialization}
Similar to \textsc{FlashAttention-2}, the forward computation in \textsc{FlashAttention-3} exhibits high parallelism across the dimensions of batch size, the number of heads, as well as the query sequence length. As such, a description at the CTA level suffices for understanding the algorithm, where each CTA processes a sub-matrix $\mathbf{Q}_i$ of the query matrix, generating the corresponding sub-matrix $\mathbf{O}_i$ of the output. For clarity, we first outline the warp-specialization approach using a circular SMEM buffer, excluding the additional complexity introduced by GEMM-softmax overlapping. Given head dimension $d$, sequence length $N$, and selecting a query block size $B_r$, the query matrix $\mathbf{Q}$ is divided into $T_r = \lceil \frac{N}{B_r} \rceil$ blocks, denoted as $\mathbf{Q}_1, ..., \mathbf{Q}_{T_r}$.

\begin{algorithm}[H]
    \caption{\small\label{alg:flash3_wgmma_ws_only}\textsc{FlashAttention-3} forward pass \textbf{without} intra-consumer overlapping -- CTA view}
    \begin{algorithmic}[1]
\REQUIRE Input matrices $\mathbf{Q}_i \in \mathbb{R}^{B_r \times d}$, $\mathbf{K}, \mathbf{V} \in \mathbb{R}^{N \times d}$ stored in HBM; select key block size $B_c$ such that $T_c = \lceil \frac{N}{B_c} \rceil$.
\STATE Construct a pipeline controller to facilitate synchronization and manage a circular SMEM buffer with $s$ stages.
\IF {the warpgroup acts as producer}
\STATE Free up a prescribed number of registers as planned.
\STATE Initiate transfer of $\mathbf{Q}_i$ from HBM into shared memory.
\STATE After the load finishes, signal to consumer warp(s) that $\mathbf{Q}_i$ is available.
\FOR{$j$ from $0$ to $T_c-1$}
    \STATE Block until the buffer's $(j\,\%\,s)$th stage is fully processed.
    \STATE Trigger the load of $\mathbf{K}_j, \mathbf{V}_j$ from HBM into the corresponding shared memory buffer stage indexed by $(j\,\%\,s)$.
    \STATE On completion, notify consumers that $\mathbf{K}_j, \mathbf{V}_j$ have been fetched.
\ENDFOR
\ELSE \STATE Assign registers based on the determined consumer warp count.
\STATE On the chip, set $\mathbf{O}_i \in \mathbb{R}^{B_r \times d}$ to zero and $\ell_i, m_i \in \mathbb{R}^{B_r}$ to $(0), (-\infty)$, respectively.
\STATE Await availability of $\mathbf{Q}_i$ in shared memory.
\FOR{$j$ from $0$ to $T_c-1$}
\STATE Wait until $\mathbf{K}_j$ becomes available in shared memory.
\STATE Compute attention scores with $\mathbf{S}_i^{(j)} = \mathbf{Q}_i \mathbf{K}_j^T$ (SS-GEMM). Sync and continue. \label{alg:ws_only_gemm1}
\STATE Backup $m_i$ as $m_i^{\mathrm{old}}$, then update via $m_i = \mathrm{max}(m_i^{\mathrm{old}}, \mathrm{rowmax}(\mathbf{S}_i^{(j)}))$. \label{code-ws:softmax_start}
\STATE Calculate $\widetilde{\mathbf{P}}_i^{(j)} = \mathrm{exp}(\mathbf{S}_i^{(j)} - m_i)$; update $\ell_i = \mathrm{exp}(m_i^{\mathrm{old}} - m_i) \ell_i + \mathrm{rowsum}(\widetilde{\mathbf{P}}_i^{(j)})$. \label{code-ws:softmax_end}
\STATE Wait for shared memory load of $\mathbf{V}_j$ to complete.
\STATE Update outputs with $\mathbf{O}_i = \mathrm{diag}(\mathrm{exp}(m_i^{\mathrm{old}} - m_i))^{-1} \mathbf{O}_i + \widetilde{\mathbf{P}}_i^{(j)} \mathbf{V}_j$ (RS-GEMM). Sync and proceed. \label{alg:ws_only_gemm2}
\STATE Mark the $(j\,\%\,s)$th pipeline buffer slot as ready for the next producer load.
\ENDFOR
\STATE Final normalization: $\mathbf{O}_i = \mathrm{diag}(\ell_i)^{-1} \mathbf{O}_i$, and set $L_i = m_i + \log(\ell_i)$.
\STATE Commit $\mathbf{O}_i$ and $L_i$ to HBM as the $i$th segment of $\mathbf{O}$ and $L$.
\ENDIF
\end{algorithmic}
\end{algorithm}

In our adaptation of \cref{alg:flash3_wgmma_ws_only} for the Hopper architecture, (de)allocation tasks utilize \verb|setmaxnreg|, TMA is employed to load $\mathbf{Q}_i$ and $\{ \mathbf{K}_j, \mathbf{V}_j \}_{0 \leq j < T_c}$, while GEMM operations within the consumer main loop are carried out through WGMMA, where SS or RS denotes whether the primary operand resides in shared memory or the register file, respectively. To understand the runtime behavior of \cref{alg:flash3_wgmma_ws_only}, it is important to recognize that the asynchronous nature of TMA loading prevents load operations from blocking one another. Additionally, the initial $s$ steps in the producer main loop require no wait instructions since the buffer is still undergoing population.

\paragraph{Pingpong scheduling} The asynchronous characteristics of WGMMA and TMA, combined with warp specialization, provide an avenue for overlapping the execution of the softmax function by one warpgroup with the GEMM computation performed by another. To illustrate the reasoning, note that, on modern computational accelerators, matmul operations achieve far greater throughput than non-matmul operations. For instance, the H100 SXM5 GPU is capable of reaching 989 TFLOPS for FP16 matrix multiplications, yet its throughput for special functions like the exponential—required in softmax—is limited to only 3.9 TFLOPS.\footnote{According to the CUDA programming guide, each streaming multiprocessor (SM) can execute 16 special function operations per cycle. Calculating with 132 SMs operating at 1830 MHz yields the 3.9 TFLOPS value for these operations.} When performing the attention forward pass in FP16 with a head dimension of 128, the computation involves 512 times more FLOPS for matmul than for exponential operations, while the exponential function operates at only 1/256th of the throughput; this means the exponential operation can occupy as much as half the time consumed by matmul. This imbalance is exacerbated in the FP8 format: although the matmul throughput is doubled, the performance of exponential operations remains unchanged.

Because the exponential operation is handled by a distinct hardware component known as the multi-function unit, it is preferable to align the execution of the exponential computation with the period when the Tensor Cores are carrying out the matrix multiplication. To achieve this concurrency, we insert synchronization barriers (\texttt{bar.sync} instructions) to ensure that the GEMMs (specifically, GEMM1—$\mathbf{P} \mathbf{V}$ from one step, and GEMM0—$\mathbf{Q} \mathbf{K}^\top$ from the next) associated with warpgroup 1 execute before those of warpgroup 2. With this ordering, softmax operations for warpgroup 1 are processed in parallel with GEMMs for warpgroup 2. The warpgroups then alternate, so when warpgroup 2 computes softmax, warpgroup 1 is engaged in GEMMs—an approach we refer to as ``pingpong'' scheduling, as depicted in~\cref{fig:pingpong_scheduling}. While real-world execution does not achieve the perfect overlap shown in the figure, we observe notable performance gains, such as improvements from 570 TFLOPS to 620–640 TFLOPS for FP16 forward passes with a head dimension of 128 and sequence length of 8192.

\paragraph{Attention variants} In the cases of multi-query attention~\citep{shazeer2019fast} and grouped query attention~\citep{ainslie2023gqa}, our methodology aligns with that of \textsc{FlashAttention-2}, implementing modified tensor indexing to prevent redundant storage of $\mathbf{K}$ and $\mathbf{V}$ within HBM.

\subsection{Intra-warpgroup overlapping GEMMs and softmax}

It is possible to interleave certain instructions from the softmax computation with those from the GEMMs, even when operating within a single warpgroup. Here, we outline a method for achieving this overlap.

In the attention algorithm, operations within the inner loop (main loop) have sequential dependencies that impede parallelization within a single iteration. For example, (local) softmax (lines \ref{code-ws:softmax_start} to \ref{code-ws:softmax_end}) relies on the output $\mathbf{S}_i^{(j)}$ of the first GEMM, while the second GEMM takes its result $\widetilde{\mathbf{P}}_i^{(j)}$ as an operand. Indeed, the wait statements in lines \ref{alg:ws_only_gemm1} and \ref{alg:ws_only_gemm2} of \cref{alg:flash3_wgmma_ws_only} serialize the execution of softmax and GEMMs. However, we can break these dependencies by pipelining across iterations through additional buffers in registers. Pursuing this idea, we propose the following two-stage\footnote{Note that the number of stages of the overlapping scheme is bounded by, but need not equal, the number $s$ of stages in the circular SMEM buffer.} GEMM-softmax pipelining algorithm:

\begin{algorithm}[H]
    \caption{\small\label{alg:flash3_wgmma}\textsc{FlashAttention-3} consumer warpgroup forward pass}
    \begin{algorithmic}[1]
      \REQUIRE  Matrices $\mathbf{Q}_i \in \mathbb{R}^{B_r \times d}$ and $\mathbf{K}, \mathbf{V} \in \mathbb{R}^{N \times d}$ present in HBM, with block size for keys given by $B_c$, and $T_c = \lceil \frac{N}{B_c} \rceil$ blocks in total.
      \STATE Adjust register allocation according to the number of consumer warps in operation.
      \STATE Initialize on-chip memory: set $\mathbf{O}_i = (0) \in \mathbb{R}^{B_r \times d}$, $\ell_i = (0)$ and $m_i = (-\infty)$ in $\mathbb{R}^{B_r}$.
      \STATE Ensure both $\mathbf{Q}_i$ and $\mathbf{K}_0$ are available in shared memory before proceeding.
      \STATE Using WGMMA, calculate $\mathbf{S}_{\mathrm{cur}} = \mathbf{Q}_i \mathbf{K}_0^T$, then commit and synchronize.
      \STATE Release the buffer's stage $0$ designated for $\mathbf{K}$.
      \STATE Determine $m_{i}$, $\tilde{\mathbf{P}}_{\mathrm{cur}}$, and $\ell_{i}$ using $\mathbf{S}_{\mathrm{cur}}$ and update the scaling of $\mathbf{O}_i$.
      \FOR{$1 \le j < T_c - 1$}
        \STATE Wait for $\mathbf{K}_j$ to become available in shared memory.
        \label{code:mainloop_start}
        \STATE Execute WGMMA to obtain $\mathbf{S}_{\mathrm{next}} = \mathbf{Q}_i \mathbf{K}_{j}^T$; commit, but do not block.
        \label{code:first_wgmma}
        \STATE Wait until $\mathbf{V}_{j-1}$ has been loaded into shared memory.
        \STATE With WGMMA, accumulate: $\mathbf{O}_{i} = \mathbf{O}_{i} + \tilde{\mathbf{P}}_{\mathrm{cur}} \mathbf{V}_{j-1}$; commit without waiting.
        \label{code:second_wgmma}
        \STATE Wait for completion of WGMMA for $\mathbf{Q}_i \mathbf{K}_{j}^T$.
        \STATE Compute the values $m_{i}$, $\tilde{\mathbf{P}}_{\mathrm{next}}$, and $\ell_{i}$ derived from $\mathbf{S}_{\mathrm{next}}$.
        \label{code:softmax}
        \STATE Once WGMMA for $\tilde{\mathbf{P}}_{\mathrm{cur}} \mathbf{V}_{j-1}$ finishes, perform rescaling of $\mathbf{O}_i$.
        \STATE Release stage $(j\,\%\,s)$ for $\mathbf{K}$ and $(j-1\,\%\,s)$ for $\mathbf{V}$ in the buffer, respectively.
        \STATE Set $\mathbf{S}_{\mathrm{cur}}$ to the contents of $\mathbf{S}_{\mathrm{next}}$.
        \label{code:mainloop_end}
      \ENDFOR \STATE Wait until $\mathbf{V}_{T_c - 1}$ is accessible in shared memory.
      \STATE Perform the computation $\mathbf{O}_{i} = \mathbf{O}_{i} + \tilde{\mathbf{P}}_{\mathrm{last}} \mathbf{V}_{T_c - 1}$ via WGMMA, then commit and synchronize.

\STATE Epilogue: Adjust $\mathbf{O}_{i}$ according to $m_{i}$. Determine $L_{i}$ utilizing both $m_{i}$ and $\ell_{i}$.
Store $\mathbf{O}_{i}$ and $L_{i}$ in HBM, placing them as the $i$-th segments of $\mathbf{O}$ and $L$ respectively.
\end{algorithmic}
\end{algorithm}

\cref{alg:flash3_wgmma} serves as a substitute for the consumer segment found in \cref{alg:flash3_wgmma_ws_only}, thereby constituting the full \textsc{FlashAttention-3} algorithm when working with FP16 precision. Conceptually, WGMMA is used to represent asynchronous GEMM throughout. During the mainloop (spanning lines \ref{code:mainloop_start} to \ref{code:mainloop_end}), the second WGMMA call within iteration $j$ (line \ref{code:second_wgmma}) is executed concurrently with the softmax computations occurring in the subsequent iteration $j+1$ (line \ref{code:softmax}).

Although the pipelined approach shown previously promises improved theoretical throughput, several practical constraints must be addressed:

\paragraph{Compiler reordering} The pseudocode presents an ideal pipeline, yet the NVCC compiler routinely reorganizes instructions for better efficiency. Such optimizations may interfere with the intended interleaving of WGMMA and non-WGMMA instructions, causing altered execution order or reduced speedup. SASS code inspection confirms that, in practice, the generated binaries exhibit the expected degree of instruction overlap (see Section~\ref{sec:2-stage-sass}).

\paragraph{Register pressure} To achieve high throughput, it's crucial to avoid excessive register spilling. Nevertheless, the dual-stage pipeline imposes extra register requirements for holding intermediate computations and passing state between pipeline steps. In particular, each threadblock needs to retain an additional $\mathbf{S}_{\mathrm{next}}$ in its registers, incurring an overhead of $B_r \times B_c \times \text{sizeof}(\text{float})$. This elevated register demand can limit the effectiveness of enlarging block sizes—a frequently used optimization—since both approaches are register-intensive. Therefore, decisions regarding block sizing and pipeline depth should be informed by profiling and benchmarking.

\paragraph{3-stage pipelining} Building on the two-stage design, we propose a three-stage alternative that incorporates further overlap between the second WGMMA and softmax steps. This extension increases Tensor Core occupancy but also intensifies register usage, given the pipeline's extra stage. The result is a more intricate trade-off between the dimensions of computation tiles and the length of the pipeline. Full algorithmic details and performance analysis for this three-stage pipeline are presented in~\cref{sec:3-stage}.

\subsection{Low-precision with FP8}
\label{sec:algofp8}

\textbf{Efficiency: layout transformations.} Executing the forward operation of \textsc{FlashAttention-3} using FP8 precision introduces unique difficulties regarding layout compatibility that are absent in the FP16 scenario.

To begin with, the tensors $\mathbf{Q}$, $\mathbf{K}$, and $\mathbf{V}$ are conventionally stored with contiguity along the head dimension. However, complying with the k-major requirement imposed by FP8 WGMMA in the second GEMM demands that $\mathbf{V}$—specifically, its tiles when placed in SMEM—be contiguous in the sequence length axis. Because the TMA load is unable to alter which dimension is contiguous, two alternatives arise: either (1) $\mathbf{V}$ must be transposed in GMEM before computation, or (2) tiles of $\mathbf{V}$ are transposed within the kernel after they are transferred to SMEM. For option (1), implementation involves either (1a) integrating the transpose operation within the epilogue of an earlier computation step, for instance, during rotary embedding, or (1b) invoking a dedicated preprocessing transpose kernel\footnote{A highly tuned transpose kernel can approach the theoretical bandwidth performance of the hardware~\citep{colfax_cutlass_transpose_2024}.} to swap the strides associated with the sequence length and head dimensions. Nevertheless, approach (1a) poses challenges for inclusion in standard libraries, while option (1b) results in significant inefficiency in memory-limited contexts such as inference workloads.

For FP8 \textsc{FlashAttention-3}, we select the second approach. Within the kernel, the transpose operation is efficiently handled by utilizing the LDSM (\verb|ldmatrix|) and STSM (\verb|stmatrix|) instructions. These instructions coordinate a warp of threads to collectively move data from shared memory (SMEM) to register memory (RMEM) and back, operating on 128-byte chunks.\footnote{The PTX specification characterizes LDSM/STSM as facilitating the transfer of $8 \times 8$ matrices with 16-bit components \cite[\S 9.7.13.4.15-16]{ptx}. By packing two 8-bit elements per 16-bit entity, we can make use of these instructions for FP8 operations. However, the transpose-capable variants of LDSM/STSM do not partition the paired 8-bit entries, requiring additional shuffling among registers between LDSM and STSM to achieve the actual tile transposition; implementation specifics are omitted here.} The LDSM/STSM primitives are frugal in register usage, enabling execution within the producer warpgroup, while also supporting layout transpositions during memory movement. Additionally, from the second iteration onward, it is possible to schedule the transpose computation for the subsequent $\mathbf{V}$ tile concurrently with the two WGMMAs that process the previous $\mathbf{V}$ and the present $\mathbf{K}$ tile.

Secondly, it is important to note that the FP32 accumulator associated with FP8 WGMMA possesses a memory layout that does not match the format expected for operand A when stored in registers, contrary to the case with FP16. Partial representations of both layouts are illustrated in \cref{fig:rmem_accum} and \cref{fig:rmem_operand}, showing how each entry resides within per-thread registers according to the specified sequence. By employing byte permute instructions, it is possible to convert the accumulator from the first WGMMA into a configuration appropriate for input to the second WGMMA, thus aligning with the arrangement of the $\mathbf{V}$ tile produced by the transpose performed within the kernel. In particular, as seen in \cref{fig:rmem_accum}, the order is modified to
$$\{ \verb|d0 d1 d4 d5 d2 d3 d6 d7| \},$$
with this register permutation consistently applied across every block of 8 bytes. From the perspective of the $\mathbf{P}$ tile's logical organization, this adjustment leads to the permutation of its columns; for example, columns $0189$ are shifted to the initial positions. To guarantee that WGMMA generates the intended output tile, the transpose operation within the kernel can similarly be set up to emit a corresponding permutation of the rows in the $\mathbf{V}$ tile.\footnote{By carrying out the transpose operation inside the kernel, additional flexibility is gained which obviates the need for shuffle instructions to transfer register ownership between threads, a challenge previously highlighted in~\citep{colfax_fp8_flashattention_2024}.}

\textbf{Accuracy: block quantization and incoherent processing.} The FP8 (e4m3) representation utilizes only 3 bits to encode the mantissa and 4 bits for the exponent, introducing greater numerical inaccuracies compared to FP16 or BF16. Furthermore, large-scale models are often characterized by the presence of outlier values~\citep{dettmers2208llm, sun2024massive}, which significantly exceed most other values, thereby complicating quantization. In response, practitioners usually employ per-tensor scaling~\citep{micikevicius2022fp8}, assigning a single scaling factor to each tensor (such as one for $\mathbf{Q}$, one for $\mathbf{K}$, and one for $\mathbf{V}$). To address the increased numerical error when computing attention in FP8, we introduce two strategies:

\begin{enumerate}

\item \textbf{Block quantization}: Each block retains a single scalar, so for $\mathbf{Q}$, $\mathbf{K}$, and $\mathbf{V}$, we partition the respective tensor into blocks with dimensions either $B_r \times d$ or $B_c \times d$, then apply quantization independently to each block. This approach can be integrated with preceding operations such as rotary embedding without incurring additional latency, owing to the memory-bandwidth limitations of rotary embedding. Because the \textsc{FlashAttention-3} algorithm is inherently designed to process blocks, we can adjust the scaling for each block in $\mathbf{S}$ to accommodate block-level quantization without introducing further computational overhead.
\item \textbf{Incoherent processing}: To mitigate the impact of outliers, we pre-process $\mathbf{Q}$ and $\mathbf{K}$ by multiplying them with a randomly chosen orthogonal matrix $\mathbf{M}$ before quantization to FP8. Due to the orthogonality property, where $\mathbf{M} \mathbf{M}^\top = I$, the computation $(\mathbf{Q} \mathbf{M}) (\mathbf{K} \mathbf{M})^\top$ is equivalent to $\mathbf{Q} \mathbf{K}^\top$, guaranteeing unchanged attention output. This strategy distributes the influence of outliers by making each component of $\mathbf{Q} \mathbf{M}$ or $\mathbf{K} \mathbf{M}$ a random linear combination of the original entries, which lowers quantization error. In implementation, following the methodology in \citet{chee2024quip} and~\citet{tseng2024quip}, we construct $\mathbf{M}$ as a product of randomly assigned diagonal matrices with $\pm 1$ and a Hadamard matrix; this allows for multiplication in $O(d \log d)$ time, rather than $O(d^2)$, and can also be combined with rotary embedding without imposing any computational penalty.
\end{enumerate}

Empirical results, as shown in \cref{sec:numerical_error}, confirm that these strategies decrease numerical errors by as much as 2.6$\times$.

\section{Empirical Validation}
\label{sec:exp}

To construct \textsc{FlashAttention-3}, we utilize CUTLASS~\citep{Thakkar_CUTLASS_2023} primitives, including the WGMMA and TMA abstractions, enabling us to assess both its performance and precision.

\begin{itemize}

\item \textbf{Benchmarking attention.} We evaluate the performance of \textsc{FlashAttention-3} by recording its runtime over a range of sequence lengths and contrasting these results with those from a baseline PyTorch implementation, \textsc{FlashAttention-2}, and a Triton-based version of \textsc{FlashAttention-2} that leverages H100-specific instructions. Additionally, comparisons are made with the cuDNN vendor-optimized \textsc{FlashAttention-2} for H100 GPUs. Our findings indicate that \textsc{FlashAttention-3} delivers up to a 2.0$\times$ speedup over \textsc{FlashAttention-2} and is 1.5$\times$ faster relative to the Triton variant. Furthermore, \textsc{FlashAttention-3} achieves peak throughput up to 740 TFLOPs/s, reaching 75\% of the theoretical computational ceiling of H100 GPUs.

\item \textbf{Ablation study.} We demonstrate through ablation analysis that the introduction of warp-specialization and GEMM-softmax pipelining notably enhances the efficiency of \textsc{FlashAttention-3}.

\item \textbf{Accuracy of FP8 attention.} Our experiments confirm that by employing block quantization and incoherent processing, the numerical error encountered in FP8 \textsc{FlashAttention-3} is reduced by a factor of 2.6$\times$.

\subsection{Benchmarking Attention}
\label{subsec:benchmark_attn}

We evaluate the execution time of various attention algorithms on an H100 80GB SXM5 GPU, using both causal and non-causal configurations and head dimensions of 64 and 128, with FP16 input representations. The benchmarking outcomes, presented in~\cref{fig:benchmark_attn_fwd} and~\cref{fig:benchmark_attn_bwd}, indicate that \textsc{FlashAttention-3} achieves speedups of approximately 1.5-2.0$\times$ over \textsc{FlashAttention-2} in the forward pass and 1.5-1.75$\times$ in the backward pass. When set against standard attention mechanisms, \textsc{FlashAttention-3} demonstrates acceleration by a factor of up to 3-16$\times$. Additionally, for sequences with medium to large lengths (1k or greater), \textsc{FlashAttention-3} outperforms even the cuDNN vendor library (a proprietary implementation optimized for H100 GPUs) in terms of speed.

\paragraph{Benchmark settings:} We vary the sequence length as 512, 1k, ..., 16k, and set batch size so that the total number of tokens is 16k. We set the hidden dimension to 2048, and head dimension to be either 64, 128, or 256 (i.e., 32 heads, 16 heads, or 8 heads). To calculate the FLOPs of the forward pass, we use:

\begin{equation*}
  4 \cdot \text{seqlen}^2 \cdot \text{head dimension} \cdot \text{number of heads}.
\end{equation*}

Applying causal masking reduces this count by half, as roughly only 50\% of the positions are processed. For the backward pass, the total FLOPs are obtained by scaling the forward pass FLOPs by a factor of 2.5, reflecting that the backward pass involves five matrix multiplications (due to recomputation), compared to two during the forward computation.

Additionally, we evaluate the runtime of FP8 during the forward pass with analogous configurations. The outcomes for headdim 256 are illustrated in ~\cref{fig:benchmark_attn_fp8}, while comprehensive results are provided in \cref{sec:benchmark_attn_fp8_full}.

\subsection{Ablation Study: 2-Stage Pipelining Experiments}
\label{sec:2-stage-pipelining-experiments}

We conduct ablation experiments on both the 2-stage WGMMA-softmax pipelining approach and the use of warp-specialization for non-causal FP16 \textsc{FlashAttention-3}, holding parameters fixed at $\{\text{batch}, \text{seqlen}, \text{nheads}, \text{hdim}\} = \{ 4, 8448, 16, 128\}$. The outcomes reported in~\cref{table:ablation_pipelining} indicate that the algorithmic enhancements—specifically the combination of asynchrony via warp-specialization and the concurrent execution of GEMM and softmax—deliver a substantial performance increase, boosting throughput from 570 to 661 TFLOPs.

\subsection{Numerical Error Validation}
\label{sec:numerical_error}

Given the ongoing discussion regarding the numerical inaccuracies in \textsc{FlashAttention}~\citep{golden2024flash}, this work evaluates \textsc{FlashAttention-2}, \textsc{FlashAttention-3}, and a conventional attention implementation by benchmarking them against a FP64 reference. In order to mimic the presence of outlier features and activations found in large language models~\citep{dettmers2208llm, sun2024massive}, the elements of $\mathbf{Q}, \mathbf{K}, \mathbf{V}$ are sampled using the following distribution:

\begin{equation*}
  \mathcal{N}(0, 1) + \mathcal{N}(0, 100) \cdot \mathrm{Bernoulli}(0.001).
\end{equation*}

Specifically, each matrix element follows a normal distribution centered at zero with a standard deviation of 1; however, for 0.1\% of the entries, we superimpose an additional independent normal noise term with a standard deviation of 10. We assess the root mean squared error (RMSE) as reported in~\cref{table:numerical_error}. For the FP16 setting, both \textsc{FlashAttention-2} and \textsc{FlashAttention-3} demonstrate RMSE values that are 1.7$\times$ lower than those from the conventional implementation, attributed to the preservation of intermediate (softmax) computations in FP32 precision. The FP8 baseline attention method utilizes per-tensor scaling, accumulates matrix multiplication results in FP32, and maintains intermediate softmax values in FP16. Owing to block quantization and less synchronized computation, \textsc{FlashAttention-3} in FP8 shows a 2.6$\times$ improvement in accuracy over the baseline.

\section{Dicussion, Limitations, Conclusion}
\label{sec:discussion}

By leveraging new programming strategies and advanced hardware capabilities—specifically asynchrony and low-precision computation—\textsc{FlashAttention-3} significantly advances both performance and accuracy in attention mechanisms. This iteration achieves a 1.5–2.0$\times$ speedup over \textsc{FlashAttention-2}, and delivers a 2.6$\times$ reduction in FP8 numerical error compared to conventional per-tensor quantization methods. Nonetheless, several aspects remain for future exploration: tailoring optimizations for LLM inference, incorporating a persistent kernel architecture into the FP8 kernel,\footnote{For benchmarking purposes, FP16 \textsc{FlashAttention-3} utilizes a persistent kernel alongside a load balancing approach, whereas FP8 \textsc{FlashAttention-3} lacks this feature. This contributes to FP8 \textsc{FlashAttention-3} performing suboptimally for short sequence lengths and causal masking, relative to FP8 cuDNN kernels.} and further investigating how low-precision attention operates in large-scale training settings. While this research centers on Hopper GPUs, the methods introduced here are likely transferable to other hardware accelerators. Advancements in primitive operations such as attention have the potential to enable novel long-context task applications through enhanced speed and accuracy.

\end{document}